\section{Banco de baterías}\label{sec:baterias}

Para seleccionar las baterías se compararon cinco modelos de ciclo profundo de dos tecnologías, plomo-ácido (inundada, AGM y gel) y LiFePO\textsubscript{4}, con un banco de \SI{24}{V} y una profundidad de descarga (DOD) de diseño de 70\,\%. Las Tablas~\ref{tab:opciones} y~\ref{tab:matriz} resumen las opciones y la matriz comparativa, con los datos de los fabricantes \cite{trojan_t105, trojan_sig, trojan_saes, ultracell, pq, victron_ds} y precios minoristas de Estados Unidos y Europa de noviembre de 2026. El costo por kWh útil se calculó como el precio entre la energía nominal multiplicada por la DOD, y el costo durante la vida útil como el costo por kWh útil entre los ciclos declarados.

\begin{table}[!htb]
\caption{Especificaciones de las opciones comerciales}
\label{tab:opciones}
\centering
\footnotesize
\setlength{\tabcolsep}{2.2pt}
\begin{tabular}{l l c c c c c}
\toprule
Ref. & Tecnología & $V$ [V] & $C$ [Ah] & Unid. 24\,V & Ciclos (DOD) & Masa [kg] \\
\midrule
B1 & Inundada & 6 & 225 & 4 & 1200 (50\,\%)\textsuperscript{a} & 28 \\
B2 & AGM & 12 & 179 & 2 & 1200 (100\,\%) & 57 \\
B3 & Gel\textsuperscript{b} & 12 & 206 & 2 & 1200 (50\,\%)\textsuperscript{a} & 62.2 \\
B4 & LiFePO\textsubscript{4} & 25.6 & 200 & 1 & 6000 (80\,\%) & 37.5 \\
B5 & LiFePO\textsubscript{4} & 25.6 & 200 & 1 & 2500 (80\,\%) & 39 \\
\bottomrule
\end{tabular}
\\[2pt]
{\footnotesize Capacidad a $C_{20}$ para B1 a B3. \textsuperscript{a}Dato de un documento del fabricante alojado por un distribuidor. \textsuperscript{b}Híbrida AGM y gel.}
\end{table}

\begin{table*}[!th]
\caption{Matriz comparativa de las baterías (módulo de \SI{24}{V})}
\label{tab:matriz}
\centering
\footnotesize
\setlength{\tabcolsep}{4pt}
\resizebox{\textwidth}{!}{%
\begin{tabular}{l r c c r r c r c >{\raggedright\arraybackslash}p{0.17\textwidth} r c}
\toprule
Ref. & Precio & $E_{\text{nom}}$ & DOD & \multicolumn{2}{c}{Costo por kWh útil [USD]} & Ciclos & $c_{\text{vida}}$ & Eficiencia & Mantenimiento & Masa & Garantía \\
\cmidrule(lr){5-6}
 & [USD] & [kWh] & tecnol. & DOD tecnol. & DOD 70\,\% & & [USD/kWh] & [\%] & & [kg] & [años] \\
\midrule
B1 & 900.00 & 5.40 & 50\,\% & 333 & 238 & 1200 & 0.28 & n.d. & Agua, ecualización, ventilación & 112 & 1\textsuperscript{a} \\
B2 & 1278.40 & 4.30 & 100\,\% & 298 & 425 & 1200 & 0.25 & n.d. & Ninguno; tensiones de carga AGM & 114 & 3\textsuperscript{b} \\
B3 & 807.25 & 4.94 & 50\,\% & 327 & 233 & 1200 & 0.27 & n.d. & Ninguno; tensiones de carga gel & 124 & n.d. \\
B4 & 809.99 & 5.12 & 80\,\% & 198 & 226 & 6000 & 0.03 & n.d. & Ninguno; BMS integrado & 38 & 5 \\
B5 & 2983.50 & 5.12 & 80\,\% & 728 & 832 & 2500 & 0.29 & 92 & Ninguno; requiere BMS externo de Victron & 39 & 3 \\
\bottomrule
\end{tabular}}
\\[2pt]
{\footnotesize Precio y masa del conjunto de unidades que forma \SI{24}{V}. n.d.: no declarado por el fabricante. El precio de B2 es el de oferta (\$639.20 por unidad) y no incluye flete.
\textsuperscript{a}Garantía de Trojan para América Latina, válida solo con DOD de hasta 50\,\%. \textsuperscript{b}Garantía de la línea AES; en Costa Rica debe confirmarse con el distribuidor.}
\end{table*}

Se seleccionó la batería B4, Power Queen LiFePO\textsubscript{4} de \SI{25.6}{V} y \SI{200}{Ah}. Tiene el menor costo por kWh útil del conjunto, \SI{198}{USD/kWh}, frente a 298, 327 y \SI{333}{USD/kWh} de las opciones de plomo-ácido (B2, B3 y B1) y \SI{728}{USD/kWh} de B5. También es la más barata con la DOD de 70\,\%: \SI{226}{USD/kWh} frente a 233 de B3, 238 de B1, 425 de B2 y 832 de B5. Declara 6000 ciclos al 80\,\% de DOD, frente a 1200 del plomo-ácido, por lo que su costo durante la vida útil es de \SI{0.03}{USD/kWh} con el precio de Estados Unidos y de \SI{0.08}{USD/kWh} con el de Costa Rica, frente a 0.25 a \SI{0.29}{USD/kWh} de las demás. Esta cifra la declara solo el fabricante, y B5 declara 2500 ciclos a la misma DOD. Aun con 2500 ciclos, B4 costaría \SI{0.08}{USD/kWh} con el precio de Estados Unidos y \SI{0.19}{USD/kWh} con el de Costa Rica, menos que las demás. La DOD de diseño, 70\,\%, queda dentro del 80\,\% al que se declaran sus ciclos.

Pesa \SI{7.3}{kg/kWh} nominal, frente a 21 a \SI{27}{kg/kWh} del plomo-ácido. Un banco de cuatro B4 pesa \SI{150}{kg}, frente a \SI{498}{kg} de ocho B3 y \SI{570}{kg} de diez B2. B5 es comparable, con \SI{7.6}{kg/kWh}. B4 no requiere mantenimiento, trae el BMS integrado y tiene cinco años de garantía del fabricante, que debe confirmarse con el distribuidor local.

El modelo se vende en Costa Rica en 850\,000 colones \cite{kuat}, precio que se usa en el presupuesto (Sección~\ref{sec:economia}). Hay dos condiciones que tomar en cuenta. La primera es que el banco usa las cuatro unidades en paralelo que admite el fabricante \cite{pq_manual}, de modo que no admite ampliaciones. La segunda es la temperatura: B4 se carga entre 0 y \SI{50}{\celsius} y no tiene protección de carga a baja temperatura \cite{pq}. En el sitio la temperatura mínima del último año fue de \SI{7}{\celsius} y la máxima de \SI{29.5}{\celsius}, dentro del intervalo de carga, y el banco se ubicará en un cuarto techado y ventilado junto al tablero principal, sin sol directo. El dimensionamiento del banco se presenta en la Sección~\ref{sec:balance}.

La batería seleccionada es de fosfato de hierro y litio (LiFePO\textsubscript{4}). Cada celda tiene \SI{3.2}{V} nominales, de modo que ocho celdas en serie forman la batería de \SI{25.6}{V}, que es la versión comercial de \SI{24}{V}; por eso el banco se dimensiona como de \SI{24}{V} \cite{linden}. Frente al plomo-ácido, admite descargas profundas con muchos más ciclos y no requiere mantenimiento, y su química es más estable térmicamente que la de otras baterías de ion-litio, con menor riesgo de incendio \cite{zubi2018}. Necesita un sistema de gestión (BMS) que equilibre las celdas y desconecte la batería por sobretensión, subtensión, sobrecorriente o temperatura fuera de rango, y no se debe cargar por debajo de \SI{0}{\celsius} \cite{linden}. La Tabla~\ref{tab:bateria} resume las características de la batería seleccionada.

\begin{table}[!htb]
\caption{Características de la batería Power Queen \SI{24}{V} \SI{200}{Ah} \cite{pq}}
\label{tab:bateria}
\centering
\footnotesize
\begin{tabular}{l l}
\toprule
Parámetro & Valor \\
\midrule
Química & LiFePO\textsubscript{4}, BMS de \SI{200}{A} integrado \\
Tensión nominal & \SI{25.6}{V} (8 celdas de \SI{3.2}{V}) \\
Capacidad y energía nominales & \SI{200}{Ah}, \SI{5120}{Wh} \\
Carga & CC/CV, \SI{28.8 \pm 0.4}{V} \\
Corriente de carga recomendada & \SI{40}{A} (0.2C) \\
Corriente continua máxima de carga y descarga & \SI{200}{A} \\
Potencia continua máxima & \SI{5120}{W} \\
Ciclos & 4000 (100\,\%), 6000 (80\,\%), 15\,000 (60\,\% de DOD) \\
Temperatura de carga & 0 a \SI{50}{\celsius} \\
Temperatura de descarga & $-20$ a \SI{60}{\celsius} \\
Conexión máxima & 4 en paralelo, 2 en serie \\
Resistencia interna & $\leq\SI{40}{m\ohm}$ \\
Protección & IP65 \\
Masa y dimensiones & \SI{37.5}{kg}, $520\times268\times\SI{220}{mm}$ \\
Garantía & 5 años \\
\bottomrule
\end{tabular}
\end{table}
